
A replication of \textit{"Ramp Metering for a Distant Downstream
Bottleneck Using Reinforcement Learning with Value
Function Approximation"} from Zhou
et al. (2020), published in the \textit{2020} volume of the Journal of Advanced Transportation (Article ID 8813467).

\section*{Replication Summary}

\textbf{Scope of Replication} --- 
This project replicates and evaluates the claims made by the original study of
\cite{ORIGINAL_STUDY}, \emph{Ramp Metering for a Distant Downstream Bottleneck
Using Reinforcement Learning with Value Function Approximation}. The original
article proposes a reinforcement learning (RL) based controller for ramp metering
and reports improved performance when compared with conventional feedback-based ramp
metering algorithms.

\vspace{1em}

\textbf{Methodology} --- 
As no supplementary material, source code, or numerical demand data was provided,
the full experimental setup was re-implemented from scratch in Python. The
simulation environment was built using a Cell Transmission Model (CTM), while the
RL agent was implemented and trained using PyTorch. Missing parameters were either assumed or reconstructed from the published figures and
methodological description.


\vspace{1em}

\textbf{Results} --- 
The no-control and reference PI-ALINEA scenarios could be reproduced closely. However, the reported performance advantage of the
RL controller could not be confirmed when stronger benchmark controllers were
introduced. In particular, optimised PI-ALINEA variants matched or outperformed
the RL controller. 


\vspace{1em}

\textbf{What was easy} --- 
Implementing the basic CTM dynamics and reproducing the qualitative no-control
and PI-ALINEA traffic patterns was comparatively straightforward once the main
simulation parameters had been identified.

\vspace{1em}

\textbf{What was difficult} --- 
The main difficulties were reconstructing undocumented model parameters, matching
the original demand profile, implementing a stable RL training procedure, and
comparing results against the original study, which reported mostly graphical
outputs and no quantitative performance tables.

\vspace{1em}

\textbf{Communication with the original authors} --- 
No attempt was made to contact the original authors. 
%No additional material, such
%as source code, numerical demand profiles, algorithmic parameters, or
%supplementary documentation, was available beyond the published article.
